\clearpage\begingroup\SourceLeaf{377b}\ArithmeticHead{46}{L}{STÄCKEL, GAUSS ALS GEOMETER.}
\renewcommand{\thefootnote}{\arabic{footnote})}\setcounter{footnote}{1}
\begin{center}
\textls[100]{Abschnitt II.}\par\medskip
\textbf{Geometria situs.}\par\medskip
\textbf{11.}\par\smallskip
\textls[100]{Allgemeines über die Geometria situs bei \textsc{Gauss}.}
\end{center}
\noindent\hspace*{\parindent}Von den Schriften, die \textsc{Gauss} veröffentlicht hat, bezieht sich keine unmittelbar
auf die Geometria situs, und doch hat dieser Gegenstand ihn sein
ganzes Leben hindurch beschäftigt. Aus Gesprächen mit \textsc{Gauss}, die in dessen
letzte Lebensjahre, 1847 bis 1855, fallen, berichtet \textsc{Sartorius} v. \textsc{Waltershausen}:
»Eine ausserordentliche Hoffnung setzte er auf die Ausbildung der
Geometria situs, in der weite, gänzlich unangebaute Felder sich befänden, die
durch unseren gegenwärtigen Kalkül noch so gut wie garnicht beherrscht
werden könnten« (\textsc{Sartorius}, S. 88). Eine ganz ähnliche Äusserung hatte er
aber etwa 50 Jahre früher getan. Am 12. Oktober 1802 schrieb er an \textsc{Olbers}:
».\ .\ .\ auch werde nächstens ein Werk von \textsc{Carnot}, \textit{Géométrie de position}\footnote{L. \textsc{Carnot}, \textit{Géométrie de position}, Paris 1803; ins Deutsche übersetzt von H. C. \textsc{Schumacher}, 2 Bände, Altona 1810. Unter Géométrie de position versteht jedoch \textsc{Carnot} etwas anderes als die Geometria situs, nämlich Untersuchungen, die sich auf die Anwendung negativer Zahlen in der Geometrie beziehen. Später hat man vielfach auch die projektive Geometrie als Geometrie der Lage bezeichnet und ihr die Geometrie des Masses gegenübergestellt, was ebenfalls mit der Geometria situs im Sinne von \textsc{Gauss} nichts zu tun hat.}
herauskommen, worauf ich überaus begierig bin. Dieser bisher fast ganz
brachliegende Gegenstand, über den wir nur einige Fragmente von \textsc{Euler}
und einem von mir sehr hochgeschätzten Geometer \textsc{Vandermonde} haben, muss
ein ganz neues Feld eröffnen und einen ganz eigenen, höchst interessanten
Zweig der erhabenen Grössenlehre bilden« (Br. G.-O. 1, S. 103).
\clearpage\SourceLeaf{378}\ArithmeticHead{47}{R}{GEOMETRIA SITUS.}\setcounter{footnote}{0}
\noindent\hspace*{\parindent}In der Tat hatte \textsc{Euler} die Frage behandelt, ob es möglich sei, die
sieben Brücken, die in Königsberg i. Pr. über die Pregelarme führen, hinter
einander und jede nur einmal zu überschreiten\footnote{L. \textsc{Euler}, \textit{Solutio problematis ad geometriam situs pertinentis}, Comment. acad. sc. Petrop. 8 (1736), 1741, S. 128 (vorgelegt den 26. August 1735); vgl. den Artikel \textit{Situation} von \textsc{D’Alembert}, \textit{Encyclopédie méthodique}, Abteilung Math., Bd. III, Paris 1789, S. 53.}. Er hatte ferner die grundlegende
Beziehung zwischen den Anzahlen der Ecken, Kanten und Seitenflächen
eines konvexen Polyeders entdeckt und bewiesen\footnote{L. \textsc{Euler}, \textit{Elementa doctrinae solidorum}, Novi Comment. acad. sc. Petrop. 4 (1752/3), 1758, S. 109 (gelesen Berlin, den 26. Nov. 1750); \textit{Demonstratio nonnullarum proprietatum, quibus solida hedris planis inclusa sunt praedita}, ebenda, S. 140 (vorgelegt den 6. April 1752); vgl. A. L. F. \textsc{Meister}, \textit{Commentatio de solidis geometricis}, Comment. Soc.-sc. Götting. Vol. 7 (1784/85) 1786, Comm. Math. S. 1.}. Endlich hatte er
sich mit den Rösselsprüngen auf dem Schachbrett befasst\footnote{L. \textsc{Euler}, \textit{Solution d’une question curieuse qui ne paroit soumise à aucune analyse}, Hist. de l’Acad., année 1759, Berlin 1766, Mémoires, S. 310.}. An ihn anknüpfend
hatte \textsc{Vandermonde} die mathematische Behandlung des Rösselsprunges gefördert
und sein Verfahren auf die analytische Darstellung von Geweben ausgedehnt\footnote{\textsc{Ch. A. Vandermonde}, \textit{Remarques sur les problèmes de situation}, Hist. de l’Acad. année 1771, Paris 1774, S. 566. V. sagt: »\textsc{Leibniz} promit un calcul de situation et mourut sans rien publier. C’est un sujet où tout reste à faire et qui mériterait bien qu’on s’en occupât«. — Zu nennen wären ferner noch die Abhandlungen: N. \textsc{Fergola}, \textit{Nuovo metodo da risolvere alcuni problemi di sito e posizione}, Atti dell’ Acad., Napoli 1787, S. 119; \textit{Nuove ricerche sulle risoluzioni dei problemi di sito}, ebenda, S. 157 und A. N. \textsc{Giordano}, \textit{Nuovo metodo da risolvere alcuni problemi di sito e posizione}, ebenda, S. 139.}.
\par
Es seien noch zwei Äusserungen von \textsc{Gauss} angeführt, die aus der Mitte
seiner Lebensbahn überliefert sind.
\par
Am 30. Oktober 1825 berichtet \textsc{Gauss} seinem Freunde \textsc{Schumacher}, dass
er in den Untersuchungen über die allgemeine Lehre von den krummen
Flächen Fortschritte gemacht habe, und sagt: »Man muss den Baum zu allen
seinen Wurzelfäden verfolgen, und manches davon kostet mir wochenlanges
angestrengtes Nachdenken. Vieles davon gehört sogar in die Geometria situs,
ein fast noch ganz unbearbeitetes Feld« (W. VIII, S. 400).
\par
In einer Aufzeichnung im \textit{Handbuch} 19 Be, die vom 22. Januar 1833
datiert ist, heisst es: »Von der Geometria Situs, die \textsc{Leibniz} ahnte, und in
die nur einem paar Geometern (\textsc{Euler} und \textsc{Vandermonde}) einen schwachen
Blick zu tun vergönnt war, wissen und haben wir nach anderthalbhundert
Jahren noch nicht viel mehr wie nichts. Eine Hauptaufgabe aus dem
\clearpage\SourceLeaf{379}\ArithmeticHead{48}{L}{STÄCKEL, GAUSS ALS GEOMETER.}\setcounter{footnote}{0}
\noindent \textls[100]{Grenzgebiet der Geometria Situs und der Geometria Magnitudinis} wird
die sein, die Umschlingungen zweier geschlossener oder unendlicher Linien
zu zählen« (W. V., S. 605).
\par
Bei den vorstehenden Worten hat \textsc{Gauss} wohl an den Brief gedacht, den
\textsc{Leibniz} am 8. September 1679 an \textsc{Huygens} gerichtet hatte und der damals
von \textsc{Uylenbroek} veröffentlicht worden war\footnote{J. \textsc{Uylenbroek}, \textit{Chr. Hugenii aliorumque seculi XVII. virorum celebrium exercitationes mathematicae et philosophicae}, Haag 1833, Heft 1, S. 9; im Heft 2, S. 6 ist der dem Briefe beigelegte Versuch einer geometrischen Charakteristik abgedruckt. Beides findet man wieder in \textsc{Leibniz}ens \textit{Mathematischen Schriften}, herausgegeben von C. J. \textsc{Gerhardt}, 1. Abt., Bd. 2, Berlin 1850, S. 19, 20, ferner in \textsc{Grassmanns} \textit{Gesammelten mathematischen und physikalischen Werken}, Bd. I, Teil 1, Leipzig 1894, S. 417, in den \textit{Oeuvres complètes} von \textsc{Chr. Huygens}, Bd. 8, Haag 1899, S. 216, 219 und endlich bei C. J. \textsc{Gerhardt}, \textit{Der Briefwechsel von Leibniz mit Mathematikern}, Bd. I, Berlin 1899, S. 568, 570.}. \textsc{Leibniz} schreibt dort: »Je crois
qu’il nous faut encore une autre Analyse proprement géométrique ou linéaire
qui nous exprime directement \textls[100]{situm}, comme l’Algèbre exprime \textls[100]{magnitudinem}«.
\par
Die Göttinger Gelehrten Anzeigen vom Jahre 1834 enthalten eine ausführliche
Besprechung der \textsc{Uylenbroek}schen Veröffentlichung von M. \textsc{Stern}
(seit 1829 Privatdozent der Mathematik in Göttingen), den der Essay von
\textsc{Leibniz} um so mehr interessiert hatte, »als er sich erinnert, von dem grössten
Mathematiker unserer Zeit einige Ideen über Geometrie gehört zu haben, die
mit einigen hier vorkommenden durchaus übereinstimmen« (S. 1940). Hierzu
ist jedoch zu bemerken, dass \textsc{Leibniz} weniger an die Geometria situs im Sinne
von \textsc{Gauss} »als an einen geometrischen Algorithmus denkt, der für einzelne
geometrische Probleme eher eine genuine Lösungsmethode liefert, als die Methode
der gewöhnlichen analytischen Geometrie«\footnote{M. \textsc{Dehn} und P. \textsc{Heegaard}, \textit{Analysis situs}, Encyklopädie der mathematischen Wissenschaften, Bd. III, Teil 1, S. 154.}.
\par
Bei der Geometria situs besitzen wir in den wenigen uns erhaltenen Aufzeichnungen
und überlieferten gelegentlichen Äusserungen nur die Spuren
\textls[100]{ausgedehnter} Untersuchungen, die \textsc{Gauss} angestellt hatte. Dies geht auch
daraus hervor, dass er wiederholt geplant hat, darüber etwas durch den Druck
bekannt zu machen. So schreibt \textsc{Möbius} am 2. Februar 1847 an \textsc{Gauss}: »Wie
ich von W. \textsc{Weber} gehört habe, haben Sie schon vor einigen Jahren beabsichtigt,
als Einleitung oder Vorbereitung der Theorie der elektrischen oder
\clearpage\SourceLeaf{380}\ArithmeticHead{49}{R}{GEOMETRIA SITUS.}\setcounter{footnote}{0}
\noindent magnetischen Strömungen eine Abhandlung über alle möglichen Umschlingungen
eines Fadens zu schreiben. Steht es nicht zu hoffen, dass diese Abhandlung
bald erscheinen wird? Die Erfüllung dieser Hoffnung würde mir und gewiss
auch vielen Andern sehr erwünscht sein« (Brief im \textsc{Gauss}-Archiv). \textsc{Gauss}
scheint dem Gedanken einer solchen Veröffentlichung näher getreten zu sein,
hat jedoch schliesslich davon Abstand genommen. Dies ergibt sich aus einem
Briefe an \textsc{Möbius} vom 13. August 1849, in dem er diesem zunächst für die
Übersendung einer Abhandlung über die Gestalten der Kurven dritter Ordnung
dankt und ihn auffordert, in entsprechender Weise die gestaltlichen Verhältnisse
der algebraischen Kurven zu untersuchen, die in \textsc{Gaussens} Dissertation
(1799) auftreten, und dann fortfährt: »Anderes damit \textls[100]{Verwandtes} hat
mich vielfach beschäftigt, und ich wollte erst in meiner neulich [16. Juli 1849]
in der Sozietät gehaltenen Vorlesung [\textit{Beiträge zur Theorie der algebraischen
Gleichungen}] die Darstellung der Hauptmomente jener Untersuchung als dritten
Teil bestimmen; aber ich würde zur Ausarbeitung dieser Darstellung einer
viel grösseren Musse bedurft haben, als sie mir zu Gebote gestanden hat«
(W. X\textsubscript{1}, S. 109). Eine Andeutung dieser Absicht ist wohl die Stelle in Art. 3
der \textit{Beiträge}, wo \textsc{Gauss} bemerkt, die von ihm vorzutragende Beweisführung
für den Fnndamentalsatz der Algebra »gehöre im Grunde einem höhern, von
Räumlichem unabhängigen Gebiete der allgemeinen abstrakten Grössenlehre
an, deren Gegenstand die nach der Stetigkeit zusammenhängenden Grössenkombinationen
sind, einem Gebiete, welches zur Zeit noch wenig angebaut ist,
und in welchem man sich auch nicht bewegen kann ohne eine von räumlichen
Bildern entlehnte Sprache« (W. III, S. 79). Vielleicht enthält das Bruchstück
einer Abhandlung über die Konvergenz der Reihen (W. X\textsubscript{1}, S. 407—410)
einen Teil jener Untersuchungen (vgl. S. 55).
\par
Im Folgenden wird zunächst berichtet werden, was sich unmittelbar auf
Grund der nachgelassenen Aufzeichnungen und mittelbar an der Hand von
Veröffentlichungen über andere Gegenstände über die Geometria situs bei
\textsc{Gauss} sagen lässt. Alsdann soll versucht werden, dem Einfluss nachzugehen,
den mündliche Andeutungen von \textsc{Gauss} auf die Entwicklung dieses Zweiges
der Grössenlehre gehabt haben.
\clearpage\endgroup
